CONJUNTOS
NOCIÓN
Entenderemos como conjunto a la reunión,
agrupación, agregado, clase, colección o
familia de integrantes homogéneos o
heterogéneos con posibilidades reales o
abstractas, que reciben el nombre de elemento del conjunto.

DETERMINACIÓN DE UN CONJUNTO

A. Extensión o forma tabular
Se enuncia todos los elementos válidos
para conjuntos con escasa cantidad de
elementos o para aquellos que siendo
excesivamente
numerosos
(o
hasta
infinitos) poseen una cierta ley de
formación la cual resulta evidente.
Comprensión o forma constructiva
Se enuncia a sus elementos por medio
de una propiedad o cualidad común a
ellos y que le es valida únicamente a
estos.


Ejemplos:
A.
Determinar el conjunto de las cinco
vocales
B.
Determinar
el
conjunto
de
los
números impares (+) menores que
16.

Por extensión:
A = {a, e, i, o, u}
B = {1, 3, 5, 7, 9, 11, 13, 15}
Por comprensión:

A = {x/x es una vocal}

B = {x/x es un número impar < 16}

RELACIÓN DE PERTENENCIA